\documentclass[a4paper, 12pt]{article}
\usepackage{amsmath,amsthm,amssymb}
\usepackage{graphicx}
\usepackage{color}
\usepackage{dsfont}
\usepackage{CJKutf8}
\usepackage{epstopdf}
\usepackage[hang]{subfigure}
\usepackage{natbib}
% \usepackage{url}
% \usepackage[T1]{fontenc}
% \RequirePackage[colorlinks,citecolor=blue,urlcolor=blue]{hyperref}
\usepackage{ulem}
\usepackage{color}
\usepackage{textcomp}

\newtheorem{coro}{Corollary}
\newtheorem{thm}{Theorem}
\newtheorem{lem}{Lemma}
\newtheorem{defi}{Definition}
\newtheorem{prop}{Proposition}
\newtheorem{assumption}{Assumption}
\newcounter{rmk}
\newcommand\rmk[1]{\vspace*{1mm} \par \stepcounter{rmk}{\noindent \bf Remark \thermk}. {#1}\vspace*{1mm}}

% page setting
\linespread{1.4}     %双倍行距，1为单倍行距，1.2是1.5倍行距
\setlength{\oddsidemargin}{0cm}
\setlength{\evensidemargin}{0cm}
\setlength{\textwidth}{160mm}
\setlength{\textheight}{220mm}

% editing
\definecolor{red}{RGB}{139,0,18}
\definecolor{lightred}{RGB}{186,25,31}
\definecolor{blue}{RGB}{0,56,108}
\definecolor{lightblue}{RGB}{69,100,139}
\renewcommand\emph[1]{{\color{red}\itshape #1}}
\newcommand\blue[1]{\textcolor{blue}{#1}}
\newcommand\red[1]{\textcolor{red}{#1}}

% define some notations
\def\e{\epsilon}
\def\t{ \mathrm{\scriptscriptstyle T} }
\def\lam{\lambda}
\def\mR{\mathds{R}}
\def\mN{\mathcal{N}}
\def\t{\mathrm{\scriptscriptstyle T} }
\def\argmin{\mathop{\arg\min}}
\newcommand\mA{\mathcal{A}}


\title{Assessing Stratified Effect in Randomized Experiment with Interference and Non-compliance}
\author{}
\date{August 2025}

\begin{document}

\maketitle

\section{Introduction}

Recent research has made significant advances in causal inference frameworks that address interference and non-compliance. 

Interference occurs when an individual's outcome is influenced not only by their treatment but also by the treatments of others. Even with unspecified interference, standard estimators have good large-sample properties (Sävje 2021; Fan, Leng, and Wu 2025). Under partial interference restricted to within clusters, cluster-based estimators provide better estimates of treatment effects (Basse and Feller 2018; Bargagli-Stoffi, Tortù, and Forastiere 2023; Jiang and Wang 2024).

Non-compliance refers to units not following their assigned treatment, adding complexity to estimating treatment effects. Instead of ignoring actual treatment adherence, some model-assisted methods estimate the effect on compliers under certain assumptions (Ren 2024). It is also possible to consider interference and non-compliance simultaneously (Imai, Jiang, and Malani 2021; Hoshino and Yanagi 2024).

However, these studies have not considered stratification within clusters, assuming that individuals within each cluster have similar responses to treatment. Stratification is a common design feature that improves estimation precision (Liu and Yang 2020; Bai et al. 2025; Shi and Li 2024), and its interaction with cluster-level interference and non-compliance has received little attention. In real-world applications, researchers may be more interested in examining complier average direct effects across various strata, such as gender, and their individual impacts. We extend the framework of Imai, Jiang, and Malani (2021) and Su and Ding (2021) to this stratified setting.

\section{Framework and notation}  %首单词首字母
% 每个session（每个subsection）开头首先需要介绍本章节主要工作
% 在介绍framework的时候，一定要注意，要把读者当新人

In this section, we review ITT analysis of two-stage randomized experiments, and the CADE causal quantity. Based on this framework, we introduce stratification into the two-stage randomized experiments, offering an improved estimation of CADE.


\subsection{Stratified two-stage randomized experiment}
% 解释cluster & stratification，在 cluster 后说明分配mechanism

In this subsection, we describe the setup of a stratified two-stage randomized experiment. We introduce cluster structure, stratification within clusters, and the two-stage treatment assignment mechanism.

We consider a stratified two-stage randomized experiment with $N$ units in total and $J$ clusters, where $n_j$ denotes the number of units in cluster $j$. Each cluster represents a group within which interference between units may occur, corresponding to a specific area or community in real-world applications. 



In the first stage of randomization, we randomly assign one of the treatment mechanisms to each cluster. We consider two mechanisms, indicated by $A_j \in \{0,1\}$, where $A_j = 0$ indicates cluster $j$ is assigned to a low proportion of treatment, and $A_j = 1$ for the high proportion, where $J_a$ clusters are assigned to mechanism $a$. In that case, we have $J_0 + J_1 = J$. The probability of $J$-dimensional mechanism assignment vector $\mathbf{A} = (A_1, A_2, \cdots, A_J)$ takes a particular value $\mathbf{a} = (a_1, a_2, \cdots, a_J)$ is  
% 暂时想不到好的方式替代\tbinom
% $$\Pr(\mathbf{A} = \mathbf{a}) = \frac{1}{\tbinom{J}{J_1}}$$
$$\Pr(\mathbf{A} = \mathbf{a}) = \frac{1}{\dbinom{J}{J_1}}$$

We introduce a stratification scheme consisting of $K$ strata, beyond the clustering structure. Strata are typically based on characteristics such as gender, where units in different strata may differ in treatment effects or their response to covariates. Each unit is uniquely associated with one of the $J$ clusters and simultaneously belongs to one of the $K$ strata. We denote $B_{ij} = k$ to indicate that unit $i$ in cluster $j$ belongs to stratum $k$. Let $\mathbf{B}_j = \{B_{1j}, \dots, B_{n_jj}\}$ represent the $n_j$-dimensional vector of stratification for cluster $j$. Here, $n_{j[k]}$ denotes the number of units in stratum $k$ within cluster $j$, satisfying the condition $\sum_{k=1}^K n_{j[k]} = n_j$.
% TODO
In the second stage, considering stratification, we assign the treatment status of units in each cluster.
In cluster $j$, we assume $n_{j[k]1}$ units are sampled without replacement and assigned to treatment from stratum $k$, while the other $n_{j[k]0}$ units receive the control. The total $n_{j1} = \sum^{K}_{k=1} n_{j[k]1}$ units receive treatment in cluster $j$. We use a binary variable $Z_{ij}$ to denote the treatment assignment of unit $i$ in cluster $j$, and $\mathbf{Z}_j = (Z_{1j}, Z_{2j} \dots Z_{n_jj})$ is the $n_j$ dimension vector of treatment assignment, with distribution 

% 确定n_j[k]1
$$\Pr(\mathbf{Z}_j = \mathbf{z}_j) = \prod^{K}_{k=1} \frac{n_{j[k]1}!\, n_{j[k]0}!}{n_{j[k]}!}$$

where $\mathbf{z}_j$ is a particular value of $\mathbf{Z}_j$.

% \red{Note: Now using $[$ to differ $n_{j[k]}$ from $n_{jz}$}
%所有的assumption，结论，提出之前，什么地方需要这些假设
%提出之后，这些假设的参考文献、他们的用处，等等



We use the finite population framework of causal inference. To account for non-compliance, let $D_{ij}$ denote the treatment receipt and $Y_{ij}$ the observed outcome of unit $i$ in cluster $j$. We write $\mathbf{D}_j = (D_{1j}, D_{2j}, \dots, D_{n_{j}j})$ and $\mathbf{Y}_j = (Y_{1j}, Y_{2j}, \dots, Y_{n_jj})$ for the treatment receipt and observed outcome vectors of cluster $j$, respectively. 

With the potential outcomes framework of causal inference, we have $Y_{ij}(\mathbf{z}, \mathbf{d})$ and $D_{ij}(\mathbf{z})$ representing the potential outcome and treatment receipt of unit $i$ in cluster $j$, when the treatment assignment and receipt vectors of all $N$ units equal $\mathbf{z}$ and $\mathbf{d}$, respectively. Without any restriction on interference, each unit has $2^N$ potential values of treatment receipt and outcome; we therefore assume that interference occurs only within clusters. (Imai 2017; Hong and Raudenbush 2006; Sobel 2006; Hudgens and Halloran 2008)


\begin{assumption}{\quad}


    $$
    Y_{ij}(\mathbf{z}) = Y_{ij}(\mathbf{z}')
    $$
    $$
    D_{ij}(\mathbf{z}) = D_{ij}(\mathbf{z}')
    $$
    for all $\mathbf{z}_j = \mathbf{z}_j '$.
    
We assume the treatment receipt and treatment of a unit could be only influenced by units in the same cluster, and units in other clusters has no effects on them. 
\end{assumption}

\subsection{Intention-to-treat effect with stratification}

We next review previous findings regarding the ITT analysis of two-stage randomized experiments under the partial interference assumption (Hudgens and Halloran 2008; Imai, Jiang, and Malani 2021). Under stratification, we give new stratified identifications of the ITT effects, analogous to the regression-adjusted estimators for stratified experiments studied by Liu and Yang (2020) and Su and Ding (2021).

% \subsubsection{Intention-to-treat causal quatities}
%介绍这一节要做什么
For unit $i$ in cluster $j$, let $\mathbf{Z}_{-i,j} = (Z_{lj})_{l \neq i}$ denote the $(n_j - 1)$-dimensional treatment assignment vector of all units in cluster $j$ other than unit $i$, so that
$$
D_{ij}(\mathbf{z}) = D_{ij}(Z_{ij} = z_{ij},\, \mathbf{Z}_{-i,j} = \mathbf{z}_{-i,j}).
$$
We define the average potential treatment receipt for unit $i$ in cluster $j$, when unit $i$ is assigned to treatment condition $z$ under mechanism $a$, as
% 在这里define
 
$$\bar{D}_{ij}(z, a) = \sum_{\mathbf{z}_{-i,j} \in \mathbf{Z}_{-i,j}} D_{ij}(Z_{ij} = z,\, \mathbf{Z}_{-i,j} = \mathbf{z}_{-i,j})\, \Pr(\mathbf{Z}_{-i,j} = \mathbf{z}_{-i,j} \mid Z_{ij} = z, A_j = a)
$$

Similarly, we define the average potential outcome for unit $i$ in cluster $j$ as

$$\bar{Y}_{ij}(z, a) = \sum_{\mathbf{z}_{-i,j} \in \mathbf{Z}_{-i,j}} Y_{ij}(Z_{ij} = z,\, \mathbf{Z}_{-i,j} = \mathbf{z}_{-i,j})\, \Pr(\mathbf{Z}_{-i,j} = \mathbf{z}_{-i,j} \mid Z_{ij} = z, A_j = a)
$$

Here, $z \in \{0,1\}$ denotes the scalar treatment assignment of unit $i$, and $\bar{D}_{ij}(z, a)$ and $\bar{Y}_{ij}(z, a)$ are marginalized over the assignments of all other units in the cluster under mechanism $a$. All dependence on $\mathbf{Z}_{-i,j}$ has been removed, and all subsequent stratum-, cluster-, and population-level quantities are defined as averages of these unit-level quantities, indexed by the same scalar $z$ and mechanism $a$.

Given these unit-level average values of potential outcomes, we consider the stratum-, cluster-, and population-level average potential values of treatment receipt and outcome. For each stratum $k$ in cluster $j$:

% noted $K$, 

$$
D_{j[k]}(z, a) = \frac{1}{n_{j[k]}} \sum_{i=1}^{n_j} \bar{D}_{ij}(z, a) \mathbf{1}(B_{ij} = k),
$$

$$
Y_{j[k]}(z, a) = \frac{1}{n_{j[k]}} \sum_{i=1}^{n_j} \bar{Y}_{ij}(z, a) \mathbf{1}(B_{ij} = k),
$$

To obtain the cluster-level and population-level averages, we weight each stratum by its size within the cluster:

$$
D_j(z, a) = \frac{1}{n_j} \sum_{k=1}^K n_{j[k]} D_{j[k]}(z, a),
\quad
Y_j(z, a) = \frac{1}{n_j} \sum_{k=1}^K n_{j[k]} Y_{j[k]}(z, a),
$$

Finally, the population-level averages are calculated by aggregating across all clusters, which can also be expressed in terms of stratified statistics:

% todo
$$
D(z, a) = \frac{1}{N} \sum_{j=1}^{J} n_j D_j(z, a) = \frac{1}{N} \sum_{j=1}^{J} \sum_{k=1}^K n_{j[k]} D_{j[k]}(z, a),
$$

$$
Y(z, a) = \frac{1}{N} \sum_{j=1}^{J} n_j Y_j(z, a) = \frac{1}{N} \sum_{j=1}^{J} \sum_{k=1}^K n_{j[k]} Y_{j[k]}(z, a).
$$


Following works of Hudgens and Halloran (2008), and Imai (2021), we have the ITT effects representing the average direct effect of treatment assignment on treatment receipt and outcome under mechanism $a$ as,
$$
\mathrm{DED}_{ij}(a) = \bar{D}_{ij}(1, a) - \bar{D}_{ij}(0, a),
\quad
\mathrm{DEY}_{ij}(a) = \bar{Y}_{ij}(1, a) - \bar{Y}_{ij}(0, a).
$$
where DED and DEY are the average direct effect of $D$ and $Y$ respectively. Then, we average $\text{DED}_{ij}$ and $\text{DEY}_{ij}$ to stratum and cluster level as,

% 替换 DED_ij
$$
\mathrm{DED}_{j[k]}(a) = \frac{1}{n_{j[k]}} \sum_{i=1}^{n_j} \mathrm{DED}_{ij}(a) \mathbf{1}(B_{ij} = k),
\quad
\mathrm{DED}_{j} (a) = \frac{1}{n_j} \sum_{k=1}^{K} n_{j[k]} \mathrm{DED}_{j[k]}(a).
$$

$$
\mathrm{DEY}_{j[k]}(a) = \frac{1}{n_{j[k]}} \sum_{i=1}^{n_j} \mathrm{DEY}_{ij}(a) \mathbf{1}(B_{ij} = k),
\quad
\mathrm{DEY}_{j} (a) = \frac{1}{n_j} \sum_{k=1}^{K} n_{j[k]} \mathrm{DEY}_{j[k]}(a).
$$

and $\mathrm{DED}_{j[k]}(a)$ and $\mathrm{DEY}_{j[k]}(a)$ indicate how the treatment assignment affect the outcome or receipt respectively in a specific stratum $k$ in cluster $j$ under assignment mechanism $a$.
To estimate $\text{DED}_{j[k]}(a)$ and $\text{DEY}_{j[k]}(a)$, we have natural estimators as $\widehat{\text{DED}}_{j[k]}(a) = \widehat{D}_{j[k]}(1, a) - \widehat{D}_{j[k]}(0,a)$ and $\widehat{\text{DEY}}_{j[k]}(a) = \widehat{Y}_{j[k]}(1, a) - \widehat{Y}_{j[k]}(0,a)$. By averaging the strata in cluster $j$, the stratified ITT estimator for treatment receipt $\text{DED}_{j}(a)$ and outcome $\text{DEY}_{j}(a)$ are
$$
\widehat{\text{DED}}_{j}(a) = \frac{1}{n_j} \sum_{k=1}^{K}n_{j[k]} \widehat{\text{DED}}_{j[k]}(a)
,\quad
\widehat{\text{DEY}}_{j}(a) = \frac{1}{n_j} \sum_{k=1}^{K}n_{j[k]} \widehat{\text{DEY}}_{j[k]}(a)
$$

Aggregating across the $J_a$ clusters assigned to mechanism $a$, we define the population-level ITT estimators:
$$
\widehat{\mathrm{DED}}(a) = \frac{1}{J_a}\sum_{j:\,A_j=a}\widehat{\mathrm{DED}}_j(a),
\qquad
\widehat{\mathrm{DEY}}(a) = \frac{1}{J_a}\sum_{j:\,A_j=a}\widehat{\mathrm{DEY}}_j(a).
$$

We propose the stratified CADE estimator as the ratio of these two population-level ITT estimators:
$$
\widehat{\mathrm{CADE}}(a) = \frac{\widehat{\mathrm{DEY}}(a)}{\widehat{\mathrm{DED}}(a)}.
$$

\begin{thm}
Under Assumption~1, for each cluster $j$ and mechanism $a \in \{0,1\}$, the stratified ITT estimators $\widehat{\mathrm{DED}}_j(a)$ and $\widehat{\mathrm{DEY}}_j(a)$ are unbiased for the corresponding finite-population causal quantities:
$$
\mathbb{E}\!\left[\widehat{\mathrm{DED}}_{j}(a)\right] = \mathrm{DED}_{j}(a),
\quad
\mathbb{E}\!\left[\widehat{\mathrm{DEY}}_{j}(a)\right] = \mathrm{DEY}_{j}(a).
$$
The proof is provided in Appendix~\ref{app:proof-unbiased}.
\end{thm}

\subsection{Complier average direct effect}
 Imai (2021) initially introduced the concept of Complier Average Direct Effect (CADE) and estimated it using overall Intent-to-Treat (ITT) estimators. However, when units are associated with both clusters and strata, applying this method becomes challenging. We propose the stratified CADE estimator $\widehat{\mathrm{CADE}}(a) = \widehat{\mathrm{DEY}}(a)/\widehat{\mathrm{DED}}(a)$ as a direct ratio of the population-level ITT estimators defined above, and prove its equivalence to stratified 2SLS and its consistency.

Following Imai (2021), the compliance status of unit $i$ in cluster $j$ could be expressed as 
$$
C_{ij}(\mathbf{z}_{-i,j}) = \mathbf{1}\!\left(D_{ij}(1, \mathbf{z}_{-i,j}) = 1,\, D_{ij}(0, \mathbf{z}_{-i,j}) = 0\right)$$
and the general measure of compliance of a unit is
$$
\bar{C}_{ij}(a) = \sum_{\mathbf{z}_{-i,j}\in \mathbf{Z}_{-i,j}} C_{ij}(\mathbf{z}_{-i,j}) \Pr(\mathbf{Z}_{-i,j} = \mathbf{z}_{-i,j} \mid A_j = a),
$$
for $a = 0,1$. Let $\text{CDY}_{ij}(a) = \bar{C}_{ij}(a) \cdot \mathrm{DEY}_{ij}(a)$ denote the compliance-weighted direct effect on outcome for unit $i$ in cluster $j$ under mechanism $a$. Finally we have CADE as
$$
\text{CADE}(a) = \frac{\sum_{j=1}^{J}\sum_{i=1}^{n_j}\text{CDY}_{ij}(a)}{\sum_{j=1}^{J}\sum_{i=1}^{n_j}\bar{C}_{ij}(a)}
$$

To identify CADE, we impose the following assumptions.

\begin{assumption}[Exclusion Restriction]\quad

The potential outcome of unit $i$ in cluster $j$ depends only on its own treatment receipt $d_{ij}$, not directly on the treatment assignment vector $\mathbf{z}_j$ nor on the treatment receipts of other units in the cluster:
$$Y_{ij}(\mathbf{z}_j, \mathbf{d}_j) = Y_{ij}(\mathbf{z}_j', \mathbf{d}_j)
\quad \text{for all } \mathbf{z}_j, \mathbf{z}_j',$$
$$Y_{ij}(d_{ij}, \mathbf{d}_{-i,j}) = Y_{ij}(d_{ij}, \mathbf{d}_{-i,j}')
\quad \text{for all } \mathbf{d}_{-i,j}.$$
\end{assumption}

\begin{assumption}[Monotonicity]\quad

For every unit $i$ in cluster $j$ and every realization of $\mathbf{z}_{-i,j}$:
$$D_{ij}(1,\, \mathbf{z}_{-i,j}) \geq D_{ij}(0,\, \mathbf{z}_{-i,j}).$$
No unit receives treatment when assigned to control but refuses when assigned to treatment, i.e., there are no defiers.
\end{assumption}

\begin{assumption}[Non-zero First Stage]\quad

For each cluster $j$ and mechanism $a \in \{0,1\}$:
$$\sum_{i=1}^{n_j} \bar{C}_{ij}(a) > 0,$$
i.e., each cluster contains at least one complier under each mechanism, so that $\mathrm{DED}(a) \neq 0$ and $\widehat{\mathrm{CADE}}(a)$ is well defined.
\end{assumption}

We further establish consistency of the proposed estimator under the following assumption.

\begin{assumption}[Asymptotic Independence of Stratified Randomization]\quad

As the within-stratum cluster size grows, conditioning on the treatment assignment of a single unit becomes asymptotically negligible for the conditional distribution of all other units' assignments. Formally, for every unit $i$ in cluster $j$, treatment condition $z_{ij} \in \{0,1\}$, and mechanism $a \in \{0,1\}$:
$$
\lim_{n_{j[k]}\to \infty} \left\{\Pr\!\left(\mathbf{Z}_{-i,j} = \mathbf{z}_{-i,j} \mid Z_{ij} = z_{ij},\, A_j=a\right) - \Pr\!\left(\mathbf{Z}_{-i,j} = \mathbf{z}_{-i,j} \mid A_j=a\right)\right\} = 0.
$$
\end{assumption}

\begin{thm}
Under Assumptions~1--4, as $\min_{j,k} n_{j[k]} \to \infty$, the stratified CADE estimator is consistent:
$$
\widehat{\mathrm{CADE}}(a) \xrightarrow{p} \mathrm{CADE}(a).
$$
\end{thm}

\subsection{Connections to Two-Stage Least Squares Regression}

In this section, we establish the direct connection between the proposed ratio estimator $\widehat{\mathrm{CADE}}(a)$ and two-stage least squares (2SLS), following the approach of Basse and Feller (2018) and Imai, Jiang, and Malani (2021). We extend their results to the stratified setting, where units are cross-classified by both cluster and stratum, building on regression-based analyses for stratified and clustered experiments by Liu and Yang (2020) and Su and Ding (2021).



To account for heterogeneous stratum sizes across clusters while preserving equal stratum representation, we transform the treatment receipt and outcome variables as
$$
D^*_{ij} = \frac{n_{j[k]}\,J}{n_k}\,D_{ij},
\qquad
Y^*_{ij} = \frac{n_{j[k]}\,J}{n_k}\,Y_{ij},
$$
where $n_{j[k]}$ is the number of units in stratum $k$ of cluster $j$, $n_k = \sum_{j=1}^J n_{j[k]}$ is the total number of units in stratum $k$ across all clusters, and $J$ is the total number of clusters. This differs from the unstratified transformation $D^*_{ij} = n_j J D_{ij}/N$ used in Imai (2021): instead of weighting by cluster size, we weight by the ratio of within-cluster to population stratum size, so that stratum $k$ contributes equally to the estimator regardless of how it is distributed across clusters. Note that when all strata have equal sizes ($n_{j[k]} = n_j/K$ for all $k$), the stratified and unstratified transformations coincide.



Let $\pi_{ij,k} = \mathbf{1}(B_{ij}=k) - n_k/N$ denote the centred stratum indicator for unit $i$ in cluster $j$, for $k = 1,\ldots,K$. Because $\sum_k \pi_{ij,k} = 0$, we drop the first stratum to avoid multicollinearity and retain $\pi_{ij,k}$ for $k = 2,\ldots,K$. We consider the following linear model for the adjusted treatment receipt:
\begin{align}
D^*_{ij} = \sum_{a=0,1} \Bigl\{\beta^D_{0a}
           + \beta^D_{1a}\,Z_{ij}
           + \sum_{k=2}^{K}\Bigl[
              \delta_{ka}\,\pi_{ij,k}
            + \lambda_{ka}\,Z_{ij}\,\pi_{ij,k}
            \Bigr]  \Bigr\}\mathbf{1}(A_j = a) + \xi_{ij}, \label{eq:first-stage}
\end{align}
where $\xi_{ij}$ is an error term. The centred stratum indicators absorb systematic within-cluster heterogeneity, so that the coefficient $\beta^D_{1a}$ captures the net effect of treatment assignment on receipt, averaged over strata. Fitting (\ref{eq:first-stage}) by weighted least squares with Horvitz--Thompson weights $W_{ij} = 1/(J_a \cdot n_{j,Z_{ij}})$, where $J_a$ is the number of clusters assigned to mechanism $a$ and $n_{j,Z_{ij}}$ is the number of units with the same assignment as unit $i$ in cluster $j$, yields
$$
\hat{\beta}^D_{1a} = \widehat{\mathrm{DED}}(a), \qquad a \in \{0,1\}.
$$



The corresponding model for the adjusted outcome is
\begin{align}
Y^*_{ij} = \sum_{a=0,1} \Bigl\{\beta^Y_{0a}
           + \beta^Y_{1a}\,D^*_{ij}
           + \sum_{k=2}^{K}\Bigl[
              \mu_{ka}\,\pi_{ij,k}
            + \nu_{ka}\,Z_{ij}\,\pi_{ij,k}
            \Bigr]  \Bigr\}\mathbf{1}(A_j = a) + \varepsilon_{ij}. \label{eq:second-stage}
\end{align}
Because $D^*_{ij}$ is endogenous, we use $Z_{ij}\,\mathbf{1}(A_j=a)$ as instruments for $D^*_{ij}\,\mathbf{1}(A_j=a)$ in a 2SLS procedure, replacing $D^*_{ij}$ with its fitted value $\hat{D}^*_{ij}$ from the first stage. The 2SLS coefficient $\hat{\beta}^Y_{1a}$ in model (\ref{eq:second-stage}) is the coefficient of primary interest, implemented as the 3rd (for $a=1$) and 4th (for $a=0$) coefficients of the second-stage regression in the modified \texttt{experiment} package.

\begin{thm}
Under Assumption~1, the stratified CADE estimator $\widehat{\mathrm{CADE}}(a)$ equals the 2SLS coefficient $\hat{\beta}^{Y,\,\mathrm{2SLS}}_{1a}$ in model~(\ref{eq:second-stage}):
$$
\hat{\beta}^{Y,\,\mathrm{2SLS}}_{1a} = \widehat{\mathrm{CADE}}(a) = \frac{\widehat{\mathrm{DEY}}(a)}{\widehat{\mathrm{DED}}(a)}.
$$
The proof is provided in Appendix~\ref{app:proof-2sls}.
\end{thm}

\section{Simulation study}

In this section, we conduct a simulation study to evaluate the finite-sample performance of the stratified CADE estimator in two-stage randomized experiments with non-compliance. We modify the R package \texttt{experiment} (Imai, Jiang, and Malani 2021) to estimate CADE in stratified two-stage experiments via two-stage least squares regression, and the original unstratified package serves as the comparison benchmark.

We generate the potential outcomes and binary treatment receipts according to the model:

$$Y_{ij}(d, a) = \beta_{ka} \cdot \mathbf{1}(d = 1) + \varepsilon_{ij}$$

where $\varepsilon_{ij} \sim N(0,1)$, $\beta_{k0} \sim t_3$, and $\beta_{k1} \sim \beta_{k0} + t_3$, where $t_3$ denotes the $t$-distribution with three degrees of freedom. We generate a continuous latent treatment receipt value:

$$D'_{ij}(z, a) = \gamma_{ka} \cdot \mathbf{1}(z = 1) + \eta_{ij}$$

where $\eta_{ij} \sim N(0, 0.01)$, $\gamma_{k0} \sim |t_3| \times 0.5$, and $\gamma_{k1} \sim \gamma_{k0} + |t_3| \times 0.5$, so that the monotonicity assumption is satisfied, and the binary treatment receipt is $D_{ij} = \mathbf{1}(D'_{ij} \geq 0)$. The number of clusters $J$, number of strata $K$, and within-stratum cluster size $n_{jk}$ are varied in three scenarios.

\medskip
\noindent (1) \textbf{Scenario 1.} There are many small strata. We set the number of strata to $K = 25, 50, 100$, fix the number of clusters to $J = 10$ and the within-stratum cluster size to $n_{jk} = 10$, yielding total sample sizes of $N = 2{,}500$, $5{,}000$, and $10{,}000$.

\medskip
\noindent (2) \textbf{Scenario 2.} There are two large strata. We fix the number of clusters to $J = 10$ and the number of strata to $K = 2$, and vary the within-stratum cluster size $n_{jk} = 100, 200, 500$, yielding total sample sizes of $N = 2{,}000$, $4{,}000$, and $10{,}000$.

\medskip
\noindent (3) \textbf{Scenario 3.} There are moderately many strata of moderate size. We fix the number of strata to $K = 4$ and the within-stratum cluster size to $n_{jk} = 20$, and vary the number of clusters $J = 20, 50, 100$, yielding total sample sizes of $N = 1{,}600$, $4{,}000$, and $8{,}000$.

\medskip
Both the potential outcomes and treatment receipts are generated once and then held fixed. Conditional on the fixed population, we repeatedly simulate a stratified two-stage randomized experiment for $10{,}000$ replications. In each replication, treatment assignment is conducted independently within each cluster, assigning $70\%$ of units to treatment when $A_j = 1$ and $30\%$ when $A_j = 0$. We compare two estimators: the stratified CADE estimator $\widehat{\mathrm{CADE}}_{\mathrm{strat}}(a)$ proposed in this paper, and the unstratified estimator $\widehat{\mathrm{CADE}}_{\mathrm{group}}(a)$ from the original \texttt{experiment} package. Their performance is evaluated in terms of absolute bias, standard deviation, root mean squared error, empirical coverage probability, and mean length of nominal $95\%$ confidence intervals. For each mechanism $a \in \{0, 1\}$, the bias is defined as $\widehat{\mathrm{CADE}}(a) - \mathrm{CADE}(a)$, and confidence intervals are constructed as $\widehat{\mathrm{CADE}}(a) \pm 1.96\,\widehat{\sigma}$, where $\widehat{\sigma}$ denotes the estimated standard deviation.

Each scenario is repeated for $10000$ times.
 \appendix

\section{Proof of Theorem~1}
\label{app:proof-unbiased}

It suffices to show $\mathbb{E}[\widehat{D}_{j[k]}(z,a)] = D_{j[k]}(z,a)$ for each stratum $k$ and $z\in\{0,1\}$, where the expectation is over the second-stage randomization within cluster $j$ under mechanism $a$; the results for $\widehat{\mathrm{DED}}_j(a)$ and $\widehat{\mathrm{DEY}}_j(a)$ then follow by linearity of expectation and the weighted-average structure of the estimators. We prove the case $z=1$; the case $z=0$ is symmetric.

Within stratum $k$ of cluster $j$ under mechanism $a$, exactly $n_{j[k]1}$ units are selected uniformly at random for treatment, so the marginal probability that unit $i\in\{i': B_{i'j}=k\}$ is assigned treatment is $\Pr(Z_{ij}=1)=n_{j[k]1}/n_{j[k]}$. Therefore,
$$
\mathbb{E}\!\left[\widehat{D}_{j[k]}(1,a)\right]
= \frac{1}{n_{j[k]1}}\sum_{i:\,B_{ij}=k}\mathbb{E}[D_{ij}Z_{ij}].
$$
For each such unit $i$,
\begin{align*}
\mathbb{E}[D_{ij}Z_{ij}]
&= \Pr(Z_{ij}=1)\cdot\mathbb{E}[D_{ij}\mid Z_{ij}=1]
= \frac{n_{j[k]1}}{n_{j[k]}}\,\bar{D}_{ij}(1,a),
\end{align*}
where the last equality uses the definition of $\bar{D}_{ij}(1,a)$. Substituting back,
$$
\mathbb{E}\!\left[\widehat{D}_{j[k]}(1,a)\right]
= \frac{1}{n_{j[k]1}}\cdot\frac{n_{j[k]1}}{n_{j[k]}}\sum_{i:\,B_{ij}=k}\bar{D}_{ij}(1,a)
= D_{j[k]}(1,a).
$$
An identical argument with $D$ replaced by $Y$ yields $\mathbb{E}[\widehat{Y}_{j[k]}(1,a)]=Y_{j[k]}(1,a)$. By linearity and the weighted-average definitions,
$$
\mathbb{E}[\widehat{\mathrm{DED}}_j(a)]
= \frac{1}{n_j}\sum_{k=1}^K n_{j[k]}\,\mathrm{DED}_{j[k]}(a)
= \mathrm{DED}_j(a),
$$
and similarly for $\widehat{\mathrm{DEY}}_j(a)$. \hfill$\square$

\section{Proof of Theorem~2}
\label{app:proof-2sls}

We show that the stratified 2SLS estimator from equals
$\widehat{\mathrm{CADE}}(a) = \widehat{\mathrm{DEY}}(a)/\widehat{\mathrm{DED}}(a)$.

The design matrix $X$ contains mutually exclusive columns for $A_j=1$ and $A_j=0$ (each
column is zero on the other mechanism's observations). Consequently, the weighted
2SLS normal equations for mechanism $a$ decouple from those for mechanism $1-a$, and
the 2SLS for mechanism $a$ reduces to a weighted IV regression over the $J_a$ clusters
with $A_j = a$.


Within mechanism $a$, the design matrix for the instrument block is
$[1,\; \pi_1,\ldots,\pi_{K-1},\; Z,\; Z\pi_1,\ldots,Z\pi_{K-1}]$,
where $\pi_{ij,k} = \mathbf{1}(B_{ij}=k) - n_k/N$ are the centred stratum indicators.
By the Frisch--Waugh--Lovell theorem applied to the weighted regression with weights
$W_{ij} = 1/(J_a \cdot n_{j,Z_{ij}})$, the first-stage coefficient on $Z$ equals the
WLS slope of the residual $\tilde{Z}_{ij}$ (after projecting $Z$ out of the constant and
stratum columns) on the residual $\widetilde{D}^*_{ij}$:
$$
\hat{\gamma}^{(a)} = \frac{\displaystyle\sum_{j:\,A_j=a}\sum_{i} W_{ij}\,\tilde{Z}_{ij}\,D^*_{ij}}
                          {\displaystyle\sum_{j:\,A_j=a}\sum_{i} W_{ij}\,\tilde{Z}_{ij}^2}.
$$


The adjusted treatment receipt is $D^*_{ij} = D_{ij}\cdot n_{j[k]}J/n_k$ for unit $i$
in stratum $k$ of cluster $j$.
Because the centred stratum indicators $\pi_{ij,k}$ are orthogonal to a constant within
each cluster, the projection of $Z$ onto $[1, \Pi_1,\ldots,\Pi_{K-1}]$ within cluster $j$
under mechanism $a$ yields the stratum-adjusted residual
$\tilde{Z}_{ij} = Z_{ij} - n_{j[k]1}/n_{j[k]}$,
where $n_{j[k]1}$ is the number treated in stratum $k$ of cluster $j$.
Substituting into the numerator and summing over stratum $k$ within cluster $j$:
\begin{align*}
\sum_{k=1}^K\sum_{i:\,B_{ij}=k} W_{ij}\,\tilde{Z}_{ij}\,D^*_{ij}
&= \sum_{k=1}^K \frac{1}{J_a}\cdot\frac{n_{j[k]}}{N}
   \Bigl(\widehat{D}_{j[k]}(1,a) - \widehat{D}_{j[k]}(0,a)\Bigr)\\
&= \frac{1}{J_a}\cdot\frac{1}{n_j}\sum_{k=1}^K n_{j[k]}\,\widehat{\mathrm{DED}}_{j[k]}(a)
 = \frac{1}{J_a}\,\widehat{\mathrm{DED}}_j(a).
\end{align*}
Summing over clusters and using the analogous identity for the denominator,
$$
\hat{\gamma}^{(a)} = \frac{\dfrac{1}{J_a}\displaystyle\sum_{j:\,A_j=a}\widehat{\mathrm{DED}}_j(a)}
                          {\text{(same denominator)}}
= \widehat{\mathrm{DED}}(a).
$$


The 2SLS estimator of CADE$(a)$ is the IV (Wald) ratio:
$$
\widehat{\mathrm{CADE}}^{\,\mathrm{2SLS}}(a)
= \frac{\displaystyle\sum_{j:\,A_j=a}\sum_{i} W_{ij}\,\tilde{Z}_{ij}\,Y^*_{ij}}
       {\displaystyle\sum_{j:\,A_j=a}\sum_{i} W_{ij}\,\tilde{Z}_{ij}\,D^*_{ij}}
= \frac{\widehat{\mathrm{DEY}}(a)}{\widehat{\mathrm{DED}}(a)}
= \widehat{\mathrm{CADE}}(a),
$$
where the second equality follows from applying the same algebra in Step~3 to $Y^*_{ij}$.
\hfill$\square$


\end{document}
